\section{Negative quadratic drift: annulus and buffered region}
\label{sec:direct-negative-annulus-buffer}

For \(q_2<0\), a half-orbit can pass from the secondary scale to the
primary scale through an annulus. The two endpoint rates have the same
leading term there; their comparison therefore depends on the first
radial corrections and on errors that vanish with the radial variables.
We first express these rates in the common central height. We then extend
the negative-branch estimate to its primary root and, for positive
half-orbits that cross the annulus, extend the positive-branch estimate
to larger radii. The buffered region below uses these estimates to count
zeros on the whole height interval; the root regimes of
\cref{sec:direct-negative-root-regimes} use them to compare the two
branches near positive roots and their degeneracies.

\subsection{The annulus for negative quadratic drift}
\label{sec:direct-negative-annulus}

For the rest of the secondary analysis put
\begin{equation}
\label{eq:direct-negative-secondary-scales}
 q_0=-as^3r^2,\qquad q_1=bs^2r,\qquad q_2=-s,\qquad
 0<a\le1,\quad 0\le b\le1,\quad \max\{\sqrt a,b\}=1.
\end{equation}
The positive scales \(s,r\) tend to zero independently.
For a signed radius \(\eta=\sigma s\rho\), define
\[
 u=r/\rho,\quad v=\rho,\quad
 D_\sigma(\rho)=-q(\sigma s\rho)/s^3,\quad
 K_\sigma(u,v)=1+au^2-\sigma bu-\sigma vR_q(\sigma sv).
\]
The variable \(u\) decreases as the radius moves away from the secondary
scale, whereas \(v\) increases toward the primary scale. For the uniform
estimates we regard \(u,v\) as independent base variables, including at zero.
On the curve \(u=r/\rho,\ v=\rho>0\), we have
\(K_\sigma(r/\rho,\rho)=D_\sigma(\rho)/\rho^2\).
Choose a large fixed \(K_0\) and a small fixed \(v_1>0\).
On a slightly enlarged box \(0\le u\le1/K_0\),
\(0\le v\le2v_1\), the functions \(K_\sigma\) have uniform
positive bounds and uniform derivative bounds. Each branch
entering the thin annular box therefore has positive outward
radial flux. The regional arguments below identify the central heights
in \(\mathcal H_p\) whose half-orbits reach this box. The radial Euler
derivative is
\[
 \mathcal E=\rho\partial_\rho
       =-u\partial_u+v\partial_v=:\mathcal E_{uv}.
\]
Here \(\mathcal E\) acts on the exit radius, and \(\partial_T\), with
\(T=-s\log h\), on the central height.
With \(\kappa_3=-R_q(0)=1+O(s)\), Taylor's formula gives
\begin{equation}
\label{eq:direct-annulus-drift}
 K_\sigma=1-\sigma bu+\sigma\kappa_3 v
                       +O_{\mathcal E_{uv}^j}((u+v)^2).
\end{equation}
The notation \(O_{\mathcal E_{uv}^j}\) includes every fixed
derivative order required below, with uniform constants.

\subsubsection{Transport and central height}

We begin by integrating the height growth through the annulus. The
resulting first-exit equation will relate the exit radius to the central
height and determine their order.
Put \(y=z/(s^3\rho^2)\), and use the fast time
\(d\tau=d\log\rho/s\). The extended equations are
\[
 u_\tau=-su,\qquad v_\tau=sv,\qquad
 y_\tau=(K_\sigma^{-1}-2s)y+y^2\mathcal N_\sigma,
\]
\[
 \mathcal N_\sigma=-\frac{\mathcal B_\sigma}
                  {K_\sigma(K_\sigma+\mathcal B_\sigma y)},
 \quad
 \mathcal B_\sigma=B(\sigma sv,s^3v^2y).
\]
On \(0\le y\le2\delta_{\rm an}\), with \(\delta_{\rm an}>0\) fixed small,
the linear coefficient \(K_\sigma^{-1}-2s\), the logarithmic rate
\(y_\tau/y\), and the derivative \(\partial_y y_\tau\) are bounded
below by positive constants.
We take \(\delta_{\rm an}\) also so small that \(y_0=\delta_{\rm an}\) is an
admissible thin threshold in \cref{lem:direct-adaptive-positive} for every
drift factor \(\mathcal K\) with \(1/8\le\mathcal K\le2\); this smallness
depends only on these bounds and on \(B\).
Since \(y_\tau/y=1/(K_\sigma+\mathcal B_\sigma y)-2s\),
\begin{equation}
\label{eq:direct-annulus-level-crossing}
 \frac{d\log y}{d\log\rho}
   =\frac1{s(K_\sigma+\mathcal B_\sigma y)}-2 .
\end{equation}
At \(y=\delta_{\rm an}\) the right side is positive wherever
\(0<s(K_\sigma+\mathcal B_\sigma\delta_{\rm an})<1/2\), and graphs cross
this level only upward there. This holds for both signs on the annular
box, for \(\sigma=-\) also on \(\rho\le0.9\), where \(K_-\) keeps a positive margin,
and for \(\sigma=+\) wherever \(K_+\) has a positive margin and
\(sK_+=O(s+|\eta|)\) is small.
To integrate this growth exactly, we construct a height coordinate with
uniform bounds on its required derivatives,
\[
 \Psi_\sigma=y\exp\Theta_\sigma,\qquad
 \Theta_\sigma=y\widetilde\Theta_\sigma,\qquad
 \Psi_{\sigma,\tau}=(K_\sigma^{-1}-2s)\Psi_\sigma,
 \qquad c<\Psi_{\sigma,y}<C .
\]
As in \cref{lem:direct-adaptive-positive}, extend the coefficients
by bounded smooth retractions into the enlarged positive box, preserving
the height-rate margins. The backward base flow is
\((u\e^{st},v\e^{-st})\), with order-\(j\) derivatives bounded by
\(\e^{jst}\). Thus \cref{lem:direct-contracting-height-coordinate}
applies with \(c_x=s\), after reducing \(s\), and supplies all the
stated height and inverse bounds.

Consequently, for \(\widehat z=s^3\rho^2\Psi_\sigma\),
\begin{equation}
\label{eq:direct-annulus-action}
 \widehat z(\rho)=
 \exp\left(\frac1s\int_{\rho_i}^{\rho}
                   \frac{dt}{tK_\sigma(r/t,t)}\right)
 \widehat z(\rho_i),\qquad \rho_i=K_0r.
\end{equation}
Suppose the central prefix reaches \(\rho_i\) with adapted height
\begin{equation}
\label{eq:direct-annulus-prefix}
 \widehat z(\rho_i)
       =h \e^{\mathcal A_{\sigma,i}/s}\Upsilon_\sigma(\xi_\sigma),
 \qquad
 \xi_\sigma=\e^{\mathcal A_{\sigma,i}/s}
                    h/(s^3r^2),\qquad \Upsilon_\sigma(0)=1.
\end{equation}
The positive units \(\Upsilon_\sigma\), their reciprocals, and their required height
derivatives are uniformly bounded for small \(\xi_\sigma\).
After reducing this fixed interval, the maps
\(\xi\mapsto\xi\Upsilon_\sigma(\xi)\) have uniformly positive
derivatives and uniformly bounded inverse derivatives.
On regular prefixes this preparation follows from
\cref{lem:direct-source-height}. Across a bottleneck it follows from the
composition \eqref{eq:direct-bottleneck-inverse-composition}; this case
is carried out in \cref{sec:direct-no-root-fold}. To separate the possibly
unbounded prefix action \(\mathcal A_{\sigma,i}\) from the annular
increment, write
\[
 \mathcal A_\sigma(\rho)=\mathcal A_{\sigma,i}
                +\int_{\rho_i}^{\rho}
                      \frac{dt}{tK_\sigma(r/t,t)} .
\]
At the first thin exit after \(\rho_i\), \(y=\delta_{\rm an}\),
\begin{equation}
\label{eq:direct-annulus-source}
 s\log(h/s^3)=-\mathcal A_\sigma(\rho)+2s\log\rho
       +s\log\Psi_\sigma(u,v,\delta_{\rm an})
       -s\log\Upsilon_\sigma(\xi_\sigma).
\end{equation}

Take \(u\le u_0<1/K_0\), with a sufficiently large fixed
gap from the inner cut, and write
\[
 \mathcal J_\sigma(\rho)=\mathcal A_\sigma(\rho)-\mathcal A_{\sigma,i}
       =\int_{K_0r}^{\rho}\frac{dt}{tK_\sigma(r/t,t)}.
\]
This action increment satisfies $\mathcal J_\sigma\ge c\log(\rho/(K_0r))$.
The equation for the central height becomes
\[
 \xi_\sigma\Upsilon_\sigma(\xi_\sigma)=\mathfrak b_\sigma(\rho),\qquad
 \mathfrak b_\sigma(\rho)
       :=u^{-2}\e^{-\mathcal J_\sigma(\rho)/s}\Psi_\sigma(u,v,\delta_{\rm an})>0.
\]
Since $\mathcal A_{\sigma,i}$ is constant under $\mathcal E$,
for every prescribed $j\ge1$,
\[
 \mathcal E^j\mathcal J_\sigma
       =\mathcal E_{uv}^{j-1}(K_\sigma^{-1}),\qquad
 |\mathcal E^j\mathcal J_\sigma|\le C_j.
\]
Thus $\mathcal E\log\mathfrak b_\sigma
 =2-(sK_\sigma)^{-1}+\mathcal E_{uv}\log\Psi_\sigma$,
and all higher fixed Euler derivatives of its logarithm are $O_j(s^{-1})$.
Differentiating its exponential yields
\begin{equation}
\label{eq:direct-annulus-input-jets}
 |\mathcal E^j\mathfrak b_\sigma|
       \le C_js^{-j}\mathfrak b_\sigma,\qquad
 \mathfrak b_\sigma\le Cu^{c/(2s)-2},\qquad j\ge0,
\end{equation}
after decreasing the fixed $u_0$.

The uniformly regular map
$\mathcal C_\sigma(\xi)=\xi\Upsilon_\sigma(\xi)$ fixes zero.
For small $s$, $\mathfrak b_\sigma\le1$ lies in its inverse domain;
the chain rule and \eqref{eq:direct-annulus-input-jets} give
$|\mathcal E^j\mathcal C_\sigma^{-1}(\mathfrak b_\sigma)|
\le C_js^{-j}\mathfrak b_\sigma$. Since
$\log\Upsilon_\sigma(0)=0$, the same estimate holds for its composition.
Thus, for every fixed $N,j$,
after reducing $s$,
\begin{equation}
\label{eq:direct-annulus-prefix-weight}
 |\mathcal E^j\xi_\sigma|
       +|\mathcal E^j\log\Upsilon_\sigma(\xi_\sigma)|
                         \le C_{N,j}u^N .
\end{equation}
Indeed the excess power of $u\le u_0<1$ in
\eqref{eq:direct-annulus-input-jets} absorbs $s^{-j}$.
The small expanded incoming height and the outward increase of
unscaled height show that the reconstructed prefix is thin.

The event equation gives
\[
 s\partial_{\log\rho}
       \log\bigl[s^3\rho^2\Psi_\sigma(u,v,\delta_{\rm an})\bigr]
       -1/K_\sigma
 =-(y_\tau|_{y=\delta_{\rm an}})\,
                      \partial_y\log\Psi_\sigma(u,v,\delta_{\rm an})<-c.
\]
Hence the central height decreases strictly with the exit radius, with inverse
derivatives \(O(s^j)\) in logarithmic incoming height. Scalar order
makes the first-exit domains and radius subintervals into intervals of
central heights.

Define the positive central-height rate
\(G_\sigma=-s\mathcal E\log h_\sigma\), and put
\(g_0=1-2s\). From \eqref{eq:direct-annulus-source},
\begin{equation}
\label{eq:direct-annulus-source-rate}
 \begin{split}
 G_\sigma&=K_\sigma^{-1}-2s
       -s\mathcal E_{uv}\log\Psi_\sigma
       +s\mathcal E\log\Upsilon_\sigma\\
 &=g_0+\sigma bu-\sigma\kappa_3 v
        +O_{\mathcal E^j}((u+v)^2+s(u+v)),\qquad c<G_\sigma<C .
 \end{split}
\end{equation}
For \(j\ge1\), the bound
\(\mathcal E_{uv}^j\log\Psi_\sigma=O(u+v)\) follows from
uniform base derivatives and the vanishing coefficients of
\(\mathcal E_{uv}\).

\subsubsection{The endpoint slope and its vanishing error}

The first-exit equation controls differentiation in central height.
We now follow the graph from the thin exit to the fixed section,
retaining the radial weight in the slope error. This weight is needed
when the leading terms of the two rates cancel.

\begin{lemma}[Annular endpoint and central-height rates]
\label{lem:direct-annulus-rates}
Under \eqref{eq:direct-annulus-prefix}, consider a thin exit
in the smaller annular box just specified. Let \(I_\sigma\)
be its endpoint and \(M_\sigma=-\mathcal E I_\sigma\).
Then \(M_\sigma>0\), and
\begin{equation}
\label{eq:direct-annulus-endpoint-rate}
 \mathcal E\log M_\sigma=2g_0+e_\sigma,\qquad
 |\mathcal E^j e_\sigma|
       \le C_js(1+|\log s|)(u+v),\quad 0\le j\le2,\qquad
 M_\sigma\le C(s\rho)^{2g_0} .
\end{equation}
In the common logarithmic height coordinate \(T=-s\log h\), put
\(m_\sigma=hI_{\sigma,h}=sM_\sigma/G_\sigma\) and
\(\mathfrak H_\sigma=\partial_T\log m_\sigma\).
Then
\begin{equation}
\label{eq:direct-annulus-clock-rate}
 \begin{split}
 \mathfrak H_\sigma
   &=2+\sigma(-\gamma_sbu+\delta_s\kappa_3 v)+\varepsilon_\sigma,\\
 \gamma_s&=(1-4s)/g_0^2,\qquad
 \delta_s=(3-4s)/g_0^2,\\
 |\varepsilon_\sigma|+|\partial_T\varepsilon_\sigma|
   &\le C(u+v)\{u+v+s(1+|\log s|)\}.
 \end{split}
\end{equation}
The two signs may have unrelated exit radii.
\end{lemma}

\begin{proof}
After the thin exit use \(\mathcal Y=z/|\eta|^2\) as the independent variable,
and write \(\mathfrak l=\log(|\eta|/(s\rho))\).
Until a small fixed \(\mathcal Y=m_{\rm land}\), its equation is
\begin{equation}
\label{eq:direct-annulus-landing-clock}
 \frac{d\mathfrak l}{d\mathcal Y}
 =\frac{sK_\sigma(u\e^{-\mathfrak l},v\e^{\mathfrak l})+B_\sigma \mathcal Y}
        {\mathcal Y[1-2sK_\sigma(u\e^{-\mathfrak l},v\e^{\mathfrak l})-2B_\sigma \mathcal Y]},
 \quad \mathfrak l(s\delta_{\rm an})=0,\quad
 B_\sigma=B(\sigma sv\e^{\mathfrak l},s^2v^2\e^{2\mathfrak l}\mathcal Y).
\end{equation}
The bracketed factor in the denominator has a fixed positive margin.
Integration gives
\[
 0\le\mathfrak l\le C_Bm_{\rm land}
      +C_Ks\log\frac{m_{\rm land}}{s\delta_{\rm an}}.
\]
Here \(C_B,C_K\) use only the coefficient bounds on the enlarged
drift box and the denominator margin. Choose the common landing level
so that \(C_Bm_{\rm land}<\tfrac12\log(0.95/0.9)\), then reduce
\(s\), with \(\delta_{\rm an}\) fixed, so that the second term is
below the same bound. Thus \(\mathfrak l<\log(0.95/0.9)<\log2\).
A stopped bootstrap keeps the graph in the enlarged positive annulus.
For the negative-branch extension to exits \(\rho\le0.9\) below,
use the fixed enlarged band \(v<0.95\):
\(K_-=1+au^2+bu-v+O(s)\) has a positive margin there for small
\(s\), and the stricter growth bound keeps the landing inside that
band. It lands in the
cone at
\[
 t=m_{\rm land}s^2\rho^2\e^{2\mathfrak l},\qquad U(t)=t/(2m_{\rm land}).
\]
Admissible graphs land before their endpoint: that endpoint has
\(z/|\eta|^2\) bounded below, and \(m_{\rm land}\) is chosen smaller.

At \(u=v=0\), \(K=1\) and \(B_\sigma=B_0(p)\), so both signs
have the landing shift
\[
 \mathfrak l_0=\frac{s}{g_0}\log\frac{m_{\rm land}}{s\delta_{\rm an}}
       -\frac{1}{2g_0}
          \log\frac{g_0-2B_0m_{\rm land}}{g_0-2B_0s\delta_{\rm an}}.
\]
Write \(R(u,v,\mathfrak l,\mathcal Y)\) for the right side of
\eqref{eq:direct-annulus-landing-clock}, and put
\(\epsilon_s=s(1+|\log s|)\). Every positive-order mixed
\((u,v,\mathfrak l)\) derivative of \(B_\sigma\) has a factor
\(s\), while \(K_\sigma\) occurs multiplied by \(s\).
The denominator margin therefore gives
\[
 |\partial_{u,v,\mathfrak l}^{\alpha}R|
       \le C_\alpha(s/\mathcal Y+s),\qquad 1\le|\alpha|\le4.
\]
The difference from the equation at the base origin has this bound
times \(u+v\). The finite-order variational argument used for
\eqref{eq:direct-adaptive-landing} now has fixed initial height
\(s\delta_{\rm an}\), zero initial base variations, and
\[
 \int_{s\delta_{\rm an}}^{m_{\rm land}}(s/\mathcal Y+s)\,d\mathcal Y
       \le C\epsilon_s.
\]
It gives \(\|\partial_{u,v}^{\alpha}\mathfrak l\|\le
C_\alpha\epsilon_s\) for \(1\le|\alpha|\le4\), and therefore
\(|\mathfrak l-\mathfrak l_0|\le C\epsilon_s(u+v)\).
At positive orders, \(\mathcal E_{uv}^j\) is a finite sum of
operators \(c_{i_1i_2}u^{i_1}v^{i_2}\partial_u^{i_1}\partial_v^{i_2}\), with
\(1\le i_1+i_2\le j\); each coefficient vanishes at the base
origin. Thus
\begin{equation}
\label{eq:direct-annulus-landing-weight}
 |\mathcal E_{uv}^j(\mathfrak l-\mathfrak l_0)|
       \le C_js(1+|\log s|)(u+v),\qquad 0\le j\le4.
\end{equation}
In particular \(\mathcal E\log t=2+2\mathcal E_{uv}\mathfrak l>0\).

Choose \(c_{\rm land}=1/(2m_{\rm land})>4\sup B\).
The moving-cone tail in \cref{lem:direct-weighted-tail} gives
\[
 \frac{dU(z_*;c_{\rm land}t)}{dt}
       =(z_*/t)^{2s}\tilde a_\sigma(t),\qquad c<\tilde a_\sigma<C,
\]
where, for \(1\le j\le3\),
\[
 |(t\partial_t)^j\log\tilde a_\sigma(t)|
 \le C_j(\sqrt t+q_1/\sqrt t+|q_0|/t)
 \le C_js(v+bu+u^2).
\]
At this landing the coefficients satisfy the cone smallness
condition. Hence
\begin{equation}
\label{eq:direct-annulus-exact-slope}
 M_\sigma=2(1+\mathcal E_{uv}\mathfrak l)\,
             z_*^{2s}t^{\,1-2s}\tilde a_\sigma(t).
\end{equation}
Put \(\mathfrak d=\mathcal E_{uv}\mathfrak l\), \(b_t=2(1+\mathfrak d)\), and
\(A_j=(t\partial_t)^j\log\tilde a_\sigma\). Logarithmic
differentiation of \eqref{eq:direct-annulus-exact-slope}
gives
\[
 e_\sigma=2g_0\mathfrak d+\frac{\mathcal E_{uv}\mathfrak d}{1+\mathfrak d}
          +b_tA_1,
 \qquad \mathcal E A_j=b_tA_{j+1}.
\]
Here \(\mathfrak d,\mathcal E_{uv}\mathfrak d,\mathcal E_{uv}^2\mathfrak d,
 \mathcal E_{uv}^3\mathfrak d\) are \(O(\epsilon_s(u+v))\), whereas
\(A_1,A_2,A_3=O(s(u+v))\). Two further Euler derivatives
of the displayed error involve derivatives of \(\mathfrak l\)
only through order four and of \(\log\tilde a_\sigma\) only
through order three; every term still contains one of these
vanishing factors. The denominators \(1+\mathfrak d\) stay bounded
away from zero. This proves the error bounds in
\eqref{eq:direct-annulus-endpoint-rate}; the bound on \(M_\sigma\)
follows from \eqref{eq:direct-annulus-exact-slope} and
\(t\asymp(s\rho)^2\).

To express the endpoint estimate in the common central height, use
\(\partial_T=\mathcal E/G_\sigma\). Thus
\[
 \mathfrak H_\sigma
   =\frac{\mathcal E\log M_\sigma
                      -\mathcal E\log G_\sigma}{G_\sigma}.
\]
Insert \eqref{eq:direct-annulus-source-rate} and
\eqref{eq:direct-annulus-endpoint-rate}. The constant term
is exactly \(2g_0/g_0=2\); the two linear coefficients
are those in \eqref{eq:direct-annulus-clock-rate}.
Differentiation using \(u_T=-u/G_\sigma\) and
\(v_T=v/G_\sigma\) preserves the stated weighted error.
\end{proof}

For comparable exit radii \(\omega=\Omega\rho\), with
fixed positive bounds on \(\Omega\),
\begin{equation}
\label{eq:direct-annulus-slope-calibration}
 \log\frac{M_-(\omega)}{M_+(\rho)}
 \le 2g_0\log\Omega+2C_{\rm cone}\sqrt{z_*}
               +Cs(1+|\log s|)(u+v).
\end{equation}
Indeed the shifts in \eqref{eq:direct-annulus-exact-slope}
have the same base value \(\mathfrak l_0\), and the moving-section
factor is
\[
 c_{\rm land}-\mathcal G_\sigma(c_{\rm land}t,t)
       =c_{\rm land}g_0-B_0+O(s(u+v)).
\]
The remaining factor of \(\tilde a_\sigma\) is the tail variation
\eqref{eq:direct-tail-variation} divided by \(\varpi_t\), that is, the slope
\(a_{t,p}\) of \eqref{eq:direct-centered-tail} at \(\mathcal U_\circ=c_{\rm land}\).
This landing energy lies in the common cone, so on each branch
\eqref{eq:direct-tail-first-slope} bounds the logarithm of that factor by
\(C_{\rm cone}(\sqrt{z_*}+q_1/\sqrt t)\), and \(q_1/\sqrt t\le Csu\). This
proves \eqref{eq:direct-annulus-slope-calibration}.

\begin{corollary}[Two annular exits]
\label{cor:direct-annular-pair}
Under \eqref{eq:direct-annulus-prefix} on both branches, let the thin
exits from an interval of \(\mathcal H_p\) lie in the smaller annular box
\(u_\sigma\le u_0\), \(v_\sigma\le v_0\). In \(T=-s\log h\),
\begin{equation}
\label{eq:direct-annular-pair-curvature}
 \begin{split}
 (\mathfrak H_--\mathfrak H_+)'&=
 -\frac{\gamma_sb(u_-+u_+)+\delta_s\kappa_3(v_-+v_+)}{g_0}\\
 &\quad+O\left(\sum_\sigma(u_\sigma+v_\sigma)
       \{u_\sigma+v_\sigma+s(1+|\log s|)\}\right).
 \end{split}
\end{equation}
The negative annular branch alone gives
\begin{equation}
\label{eq:direct-annular-exterior-curvature}
 |\partial_T^2[(\log\Pi_p')(I_-)]|\le C(sv_-)^{2-4s}
                    \le Cs^{3/2}v_-,\qquad s<1/8,\quad v_-<1.
\end{equation}
For any fixed \(b_{\min}>0\), choosing the annular margins small relative to
\(b_{\min}\), then reducing \(s\), gives \(\mathcal L''<0\) when \(b\ge b_{\min}\).
\end{corollary}

\begin{proof}
Differentiate \eqref{eq:direct-annulus-clock-rate} using
\(\partial_T=\mathcal E_\sigma/G_\sigma\),
\(u_{\sigma,T}=-u_\sigma/G_\sigma\), \(v_{\sigma,T}=v_\sigma/G_\sigma\),
and \(G_\sigma=g_0+O(u_\sigma+v_\sigma)\); this gives
\eqref{eq:direct-annular-pair-curvature}.
By \cref{lem:direct-annulus-rates}, \(I_{-,T}=-M_-/G_-\), and
\(|I_{-,T}|+|I_{-,TT}|\le C(sv_-)^{2-4s}\).
Equation \eqref{eq:direct-exterior-curvature} therefore gives the first
bound in \eqref{eq:direct-annular-exterior-curvature}; the second follows
from \((sv_-)^{2-4s}/v_-=s^{2-4s}v_-^{1-4s}\le s^{3/2}\).
For \(b\ge b_{\min}\), the leading negative term dominates both the weighted
error and the exterior contribution, proving \(\mathcal L''<0\).
\end{proof}

\subsubsection{The negative branch beyond the annulus}
\label{sec:negative-branch-far-root}

Beyond the annulus, negative outputs may lie in compact primary
intervals or near the primary root close to \(X=1\). We continue the
annular calculation on the compact intervals. Near the root, where the
drift tends to zero, the boundary-flow coordinate gives a negative
endpoint rate and estimates its derivative with the same drift weight
as the endpoint derivatives. These weights allow the exterior
contribution to be absorbed in the regional comparisons.

For the negative branch fix the regular primary entry at
\(\rho_{\rm pr}=0.8\), take \(0.85\) as the lower bound for the projected
far-root state \(w_n\), and use \(\rho=0.9\) as the last primary
thin-exit test. These thresholds give a fixed phase advance and overlapping
ranges of central heights.
In primary coordinates \(|\eta|=sX,\ z=s^2Y\), the negative boundary
drift is
\[
 K_n(X)=ar^2+brX+X^2+X^3R_q(-sX),\qquad F_n=K_n/X .
\]
It has a simple first far root \(X_*=1+O(r+s)\),
with \(K_n'(X_*)<0\), and is positive below it.
Construct the boundary-flow coordinate on a band from
below \(0.8\) to beyond \(X_*\), with small fixed output
height \(Y_n>0\). Its far-root interval is
\begin{equation}
\label{eq:direct-negative-far-core}
 w_n=\mathcal W_n(X(s^2Y_n;h),Y_n)\ge0.85,\qquad w_n<X_* .
\end{equation}

\begin{lemma}[Negative branch]
\label{lem:negative-branch-far-root}
Fix \(M\ge K_0\). Suppose that the negative graph of a central
height in \(\cH_p\) is thin at \(\rho=Mr\), with the prefix preparation
\eqref{eq:direct-annulus-prefix} at \(\rho_i=K_0r\).
\begin{enumerate}[label=\textup{(\roman*)}]
\item If its first thin exit occurs at \(\rho\le0.9\), then \(\rho>Mr\)
and, in \(T=-s\log h\),
\begin{equation}
\label{eq:direct-negative-primary-rates}
 \begin{split}
 \mathfrak H_-&=2-3\rho+
                O(r/\rho+s(1+|\log s|)),\\
 \mathfrak H_-'&=-3\rho(1-\rho)
                +O(r+s(1+|\log s|))
 \quad\hbox{on fixed compact primary intervals},\\
 |I_{-,T}|+|I_{-,TT}|&\le C(s\rho)^{2-4s}=o(1).
 \end{split}
\end{equation}
\item If the graph is still thin at \(\rho=0.9\), it satisfies
\eqref{eq:direct-negative-far-core} with \(w_n>0.89\). For every central
height with \eqref{eq:direct-negative-far-core},
\begin{equation}
\label{eq:direct-negative-far-rate}
 \begin{gathered}
 \mathfrak H_-=2-3w_n+O(r+s+\varepsilon_{n,s})<-c_n<0,\qquad
 \mathfrak H_-'=F_n(w_n)\{Q_n(w_n)+o(1)\},\\
 Q_n\le-c_n,\qquad |I_{-,T}|+|I_{-,TT}|\le CE_nF_n(w_n),
 \end{gathered}
\end{equation}
where \(Q_n=(XK_n''-K_n')/X^2\), \(E_n=s^2(z_*/(s^2Y_n))^{2s}=O(s^2)\),
and \(\varepsilon_{n,s}\to0\), uniformly up to the far root.
\end{enumerate}
The exit radius in (i) and the state \(w_n\) in (ii) increase with \(T\).
The constants are uniform on the chart
\eqref{eq:direct-negative-secondary-scales}.
\end{lemma}

\begin{proof}
\emph{Part (i).} Since the graph is thin at \(Mr\) and, by
\eqref{eq:direct-annulus-level-crossing}, crosses \(y=\delta_{\rm an}\) only
upward on the annulus, a first thin exit before \(\rho=0.9\) has \(\rho>Mr\).
The negative annular coordinate extends to \(\rho\le0.9\):
its drift factor has a uniform positive margin there.
Its equation for the central height gives
\[
 G_-=\frac1{1-\rho}+O_{\mathcal E^j}(r/\rho+s).
\]
The endpoint argument of
\cref{lem:direct-annulus-rates} on this enlarged compact
base box gives
\[
 \mathcal E\log M_-=2-4s+
        O_{\mathcal E^j}\bigl(s(1+|\log s|)(r/\rho+\rho)\bigr).
\]
The weighted error follows from the landing equation:
the comparison at the base origin has \(K=1\), and
the derivatives remain bounded on \(0\le\rho\le0.9\).
Inserting these expressions into
\((\mathcal E\log M_--\mathcal E\log G_-)/G_-\) gives
\eqref{eq:direct-negative-primary-rates}; the last estimate there uses
\(I_{-,T}=-M_-/G_-\) and its logarithmic rate.

\emph{Part (ii).} Every branch still thin at \(0.9\) reaches the output
\(Y=Y_n\) inside the band. Its initial projected state is
\(0.9+O(s\delta_{\rm an})>0.89\) and lies below \(X_*\).
The projected flow increases toward the root, and the inverse
projection differs by \(O(Y_n)\), remaining in the band.
The inverse-flux identity gives continuation to \(Y_n\).

At \(\rho_{\rm pr}=0.8\), write
\[
 \begin{gathered}
 \mathcal J_b=\int_{K_0r}^{\rho_{\rm pr}}\frac{dt}{tK_-(r/t,t)},
 \qquad
 \mathcal A_b=\mathcal A_{-,i}+\mathcal J_b,\\
 \xi_n=\e^{\mathcal A_b/s}\frac{h}{s^3}.
 \end{gathered}
\]
The secondary and primary expanded inputs are related by
\[
 \xi_-=\beta_b\xi_n,\qquad
 \beta_b=r^{-2}\e^{-\mathcal J_b/s}.
\]
The factor \(r^{-2}\) comes from changing the scale
\(s^3r^2\) of the central height to \(s^3\). The uniform upper bound
for the annular drift gives
\(\mathcal J_b\ge c\log(\rho_{\rm pr}/(K_0r))\). After fixing smaller
upper bounds for \(r\) and then \(s\), this implies
\[
 0<\beta_b\le r^{c/(2s)-2}\le1.
\]

For \(Y_i=z/s^2\), write
\(y_i=Y_i/s=z/s^3=\rho_{\rm pr}^2y_a(\rho_{\rm pr})\), where
\(y_a=z/(s^3\rho^2)\), and set
\[
 \mathcal C_b(y_i)
   =\rho_{\rm pr}^2\Psi_-(r/\rho_{\rm pr},\rho_{\rm pr},
                         y_i/\rho_{\rm pr}^2).
\]
This height coordinate and its inverse have uniform bounds
on the required derivatives, with \(\mathcal C_b'(0)=1\).
Annular transport gives
\[
 \mathcal C_b(y_i)
     =\xi_n\Upsilon_-(\beta_b\xi_n).
\]
As in \eqref{eq:direct-bottleneck-inverse-composition}, inversion
gives \(Y_i=s\xi_nV_n(\xi_n)\), with \(V_n(0)=1\). Since
\(\beta_b\le1\), the unit, its reciprocal, and the map inverse have
uniform derivative bounds on a fixed small \(\xi_n\)-interval.

For central heights with \eqref{eq:direct-negative-far-core}, the projected
state at \(X=\rho_{\rm pr}\) satisfies
\(w_i=\mathcal W_n(\rho_{\rm pr},Y_i)\le\rho_{\rm pr}\). The boundary
flow must advance from at most \(0.8\) to at least
\(0.85\) before the output \(Y_n\). Its bounded
speed and the projection identity give
\[
 s\log(Y_n/Y_i)\ge c_b>0,\qquad
 y_i=Y_i/s\le(Y_n/s)\e^{-c_b/s}.
\]
Thus the graph is thin at \(\rho_{\rm pr}\).
By \eqref{eq:direct-annulus-level-crossing}, it has no downward crossing
of \(y_a=\delta_{\rm an}\) on the preceding annular segment, so the whole
segment is thin. Annular transport gives
\[
 \frac{\widehat z(K_0r)}{s^3r^2}
       =\beta_b\mathcal C_b(y_i).
\]
This lies in the uniform small output domain of the
inner prefix. Its inverse recovers
\(\xi_-=O(\beta_by_i)\); uniqueness of the
flow identifies the recovered central height with \(h\). Thus
\(\xi_n=\xi_-/\beta_b\le Cy_i\) lies in the required small
domain. At the fixed output \(z=s^2Y_n\), the entry
\(X=\rho_{\rm pr}=0.8\) has \(F_n(\rho_{\rm pr})\) bounded below,
and the advance to \(w_n\ge0.85\) supplies a fixed positive phase
gap. Thus \cref{lem:direct-weighted-phase} applies with \(\delta=s\)
in the same coordinate \(T=-s\log h\). We use its drift-weighted
bound \eqref{eq:direct-phase-weight} with \(N=1\) and \(j=2\)
for the curvature estimate below; in particular,
\[
 (w_n)_T=\theta_nF_n(w_n),\qquad
 |\theta_n-1|+|\theta_n'|+|\theta_n''|
                         \le\varepsilon_{n,s}\longrightarrow0 .
\]
\Cref{lem:boundary-flow-endpoint}, with \(\ell=\delta=s\), gives
\[
 -I_{-,w_n}=E_nA_nw_nU_n(w_n),\qquad
 E_n=s^2(z_*/(s^2Y_n))^{2s}=O(s^2),\qquad
 U_n=1+O_{C^3}(s).
\]
The rate identity yields the first estimate in
\eqref{eq:direct-negative-far-rate}, uniformly up to the root, since
\(K_n'(w)/w=2-3w+O(r+s)\). By
\eqref{eq:direct-curvature-coefficient-cubic}, with \(K_n'(0)=br\) and
\(R_K(X)=R_q(-sX)\), one has
\(Q_n(X)=-br/X^2+3R_q(-sX)+O(s)<-c_n\) on the band, since
\(R_q=-1+O(s)\). The weighted phase estimate and
\eqref{eq:direct-rate-curvature} give the remaining assertions of
\eqref{eq:direct-negative-far-rate}. Finally \(\rho_T=\rho/G_->0\), and
the projected far-root state increases along the boundary flow.
\end{proof}

\subsubsection{The positive branch beyond the annulus}

A positive half-orbit may leave compact \(\rho\)-intervals and reach
unscaled radii \(\varrho\gg s\), where the drift grows cubically.
The adaptive coordinate of \cref{lem:direct-adaptive-positive} follows
this transition from the inner cut and gives the endpoint rate in both
the buffered and no-root regions, using their respective preparations
at the inner cut.

The annular action controls the expanded incoming height and its
derivatives with respect to the exit radius. This control makes the
endpoint-rate comparison uniform as the positive exit moves from the
quadratic to the cubic regime.

On the positive branch write \(\varrho=\eta=s\rho>0\), \(Q=s^3r^2\), and
\begin{equation}
\label{eq:positive-adaptive-variables}
 \varrho^+=s+\varrho,\quad u=\frac{sr}\varrho,\quad
 v=\frac\varrho{s+\varrho},\quad
 K_a=1+(1-v)(au^2-bu)-v\{R_q(\varrho)+1\}.
\end{equation}
Here \(u=r/\rho\) retains its annular meaning. In \(K_a\) and the following
lemma, \(v\) is the adaptive variable of
\eqref{eq:direct-adaptive-drift}; the annular coordinate
\(\Psi_+\) still has arguments \((r/\rho,\rho)\). With
\(\varrho_i=K_0sr\), put
\[
 \mathcal J_a(\varrho)=\int_{\varrho_i}^{\varrho}\frac{x\,dx}{-q(x)},
\]
so that the action increment of the annulus is \(\mathcal J_+(\rho)=s\mathcal J_a(s\rho)\).
At every radius the annular and adaptive normalized heights satisfy
\begin{equation}
\label{eq:positive-annular-adaptive-heights}
 \frac{z}{s\varrho^2}
   =\Bigl(1+\frac\varrho s\Bigr)\frac{z}{\varrho^2(s+\varrho)} ,
\end{equation}
so, for a half-orbit thin at \(\varrho_i\), the annular thin exit at the
same threshold \(\delta_{\rm an}\) precedes the adaptive exit of part~(iii)
below.
At an exit radius \(\varrho_e\) we write \(u_e,v_e,\varrho_e^+\) for the
values of \(u,v,\varrho^+\). The operator
\(\mathcal E=\varrho_e\partial_{\varrho_e}\) varies the exit radius at fixed
parameters, and \(T=-s\log h\).

\begin{lemma}[Positive branch beyond the annulus]
\label{lem:positive-branch-adaptive-exit}
Let \(\mathcal J\subset\cH_p\) be an interval of central heights whose
positive half-orbits are thin at \(\varrho_i\), that is,
\(z(\varrho_i;h)<\delta_{\rm an}K_0^2Q\), and satisfy
\begin{equation}
\label{eq:positive-adaptive-prefix}
 z(\varrho_i;h)=Q\,\xi\mathcal V(\xi),\qquad \xi=\e^{\mathcal A_i/s}h/Q .
\end{equation}
Here \(\mathcal A_i\) depends only on \(p\), \(\xi\) lies in a fixed
interval \([0,\bar\xi]\), \(\mathcal V(0)=1\), \(\mathcal V\) and
\(1/\mathcal V\) have uniformly bounded derivatives of the required orders
there, and \(\xi\mapsto\xi\mathcal V(\xi)\) has a uniformly positive
derivative and an inverse with uniformly bounded derivatives.
\begin{enumerate}[label=\textup{(\roman*)}]
\item One has \(-q(\varrho)=\varrho^2(s+\varrho)K_a\). On the box
\(0\le u\le3/(2K_0)\), \(0\le v\le1\), with small unscaled radius and
small \(s\), \(K_a\) lies between \(1/8\) and \(2\) and has uniform
derivative bounds, and \(K_a-1\) and its fixed Euler derivatives are
\(O(u+\varrho^+)\). Thus \cref{lem:direct-adaptive-positive} applies with
\(\mathcal K=K_a\), \(\varrho_0=sr\) and thin threshold \(y_0=\delta_{\rm an}\);
let \(\Psi_a\) be its coordinate and \(\widehat z_a\) its adapted height.
\item Put \(u_i=1/K_0\), \(v_i=K_0r/(1+K_0r)\) and \(c_i=K_0^2(1+K_0r)\),
and define \(\Upsilon_a\) by
\begin{equation}
\label{eq:positive-adaptive-unit}
 \xi\Upsilon_a(\xi)
   =c_i\Psi_a\bigl(u_i,v_i,\varrho_i,\xi\mathcal V(\xi)/c_i\bigr).
\end{equation}
Then \(\Upsilon_a(0)=1\), and \(\Upsilon_a\), its reciprocal and the
inverse of \(\xi\mapsto\xi\Upsilon_a(\xi)\) have uniformly bounded
derivatives on the interval of those \(\xi\le\bar\xi\) with
\(\xi\mathcal V(\xi)\le\delta_{\rm an}K_0^2\). Moreover
\(\widehat z_a(\varrho_i)=Q\xi\Upsilon_a(\xi)\), so that
\(\lambda_{\rm in}=\log\widehat z_a(\varrho_i)\) satisfies
\(d\lambda_{\rm in}/d\log h=1+O(\xi)\).
\item The adaptive exit, the first radius \(\varrho_e>\varrho_i\) with
\(z=\delta_{\rm an}\varrho^2(s+\varrho)\), satisfies
\begin{equation}
\label{eq:direct-no-root-adaptive-input}
 \xi\Upsilon_a(\xi)=\frac{\varrho_e^2\varrho_e^+}{Q}
      \Psi_a(u_e,v_e,\varrho_e,\delta_{\rm an})\e^{-\mathcal J_a(\varrho_e)} .
\end{equation}
Its landing height \(t=m_{\rm land}\varrho_l^2\) satisfies
\(\sqrt t+q_1/\sqrt t+|q_0|/t\le C(\varrho_e+su_e)\), and
\(|q_2|=s\); in particular \eqref{eq:direct-tail-smallness} holds there
for small \(s\).
\item Suppose in addition that \(\varrho_e\ge s\rho_{\rm g}\) on
\(\mathcal J\) and \(\mathcal J_+(\rho_{\rm g})\ge c_{\rm g}\log(1/r)\), for fixed
\(\rho_{\rm g},c_{\rm g}>0\), and that \(r\le1/2\). Then, for every
fixed \(j\) and small \(s\),
\begin{equation}
\label{eq:positive-adaptive-input-weight}
 |\mathcal E^j\xi|+|\mathcal E^j\log\Upsilon_a(\xi)|
       \le C_js^{-3-j}r^{c_{\rm g}/s-2},
\end{equation}
which is \(O(s^N)\) for every fixed \(N\). The rate
\(\mathfrak g=-\varrho_e^+\mathcal E\log h\) satisfies
\begin{equation}
\label{eq:direct-no-root-adaptive-source}
 \mathfrak g=1+O(u_e+\varrho_e^+),\qquad
 \mathcal E\log\mathfrak g=O(u_e+\varrho_e^+),\qquad
 \mathcal ET=(s/\varrho_e^+)\mathfrak g>0 .
\end{equation}
For small fixed \(m_{\rm land}\) and radius box, and after reducing
\(r\) and then \(s\), the endpoint rate
\(\mathfrak H_+=\partial_T\log m_+\) satisfies
\begin{equation}
\label{eq:direct-no-root-adaptive-rate}
 \mathfrak H_+
 =\frac{\varrho_e^+}{s\mathfrak g}
        [\,2-4s+v_e+O(\epsilon_e+u_e+\varrho_e^+)\,]>c\varrho_e^+/s>0,
\end{equation}
with \(\epsilon_e\) as in \eqref{eq:direct-adaptive-landing}.
\end{enumerate}
The constants depend on \(\mathcal V\) only through its assumed bounds.
The constants in \eqref{eq:direct-no-root-adaptive-source} and
\eqref{eq:direct-no-root-adaptive-rate}, apart from terms \(O(s^N)\), are
independent of \(\mathcal V\), \(\mathcal A_i\), \(\rho_{\rm g}\) and
\(c_{\rm g}\).
\end{lemma}

\begin{proof}
\emph{(i)} Since \((s+\varrho)(1-v)=s\) and \((s+\varrho)v=\varrho\),
\[
 \varrho^2(s+\varrho)K_a
   =s\varrho^2+as^3r^2-bs^2r\varrho-\varrho^3R_q(\varrho)=-q(\varrho).
\]
Because \(\kappa_3=-R_q(0)=1+O(s)\), one has \(R_q(\varrho)+1=O(s+\varrho)\)
with its fixed Euler derivatives; here \(\mathcal Eu=-u\),
\(\mathcal Ev=v(1-v)\) and \(\mathcal E\varrho=\varrho\). On the box,
\(|(1-v)(au^2-bu)|\le u<3/4\), since \(0<a\le1\), \(0\le b\le1\) and
\(K_0>2\). This gives the bounds on \(K_a\) and on \(K_a-1\). Along the
passage from \(\varrho_i\), \(u\) decreases from \(1/K_0\) while \(v\) and
\(\varrho\) increase, so the base point stays in the box, and
\(y_0=\delta_{\rm an}\) is admissible by the choice of \(\delta_{\rm an}\)
in \cref{sec:direct-negative-annulus}.

\emph{(ii)} At the inner section,
\(\varrho_i^2(s+\varrho_i)=Qc_i\), so the adapted height is exactly
\(Q\xi\Upsilon_a(\xi)\) by \eqref{eq:positive-adaptive-unit}.
Here \(K_0^2\le c_i\le2K_0^2\), and the stated thin interval gives
\(\xi\mathcal V(\xi)/c_i\le\delta_{\rm an}\). The uniform regularity
and positive height derivative of \(\Psi_a\), together with the assumed
bounds for \(\xi\mathcal V(\xi)\), give all the unit and inverse bounds.
The normalization \(\Psi_{a,y}(\cdot,0)=\mathcal V(0)=1\) gives
\(\Upsilon_a(0)=1\), and logarithmic differentiation gives the last claim.

\emph{(iii)} Thinness at \(\varrho_i\) allows
\cref{lem:direct-adaptive-positive} to supply the exit and landing
before the admissible endpoint. Its exact action identity, evaluated
at the exit and combined with (ii), is
\eqref{eq:direct-no-root-adaptive-input}. Since
\(\varrho_e\le\varrho_l\le2\varrho_e\),
\[
 \sqrt t\le2\varrho_e,\qquad
 q_1/\sqrt t\le Cbsu_e,\qquad |q_0|/t\le Casu_e^2.
\]
With \(u_e\le1/K_0\), these are the claimed tail bounds.

\emph{(iv)} Denote the right side of
\eqref{eq:direct-no-root-adaptive-input} by \(\mathfrak b_a\), a function
of \(\varrho_e\); through this equation \(\xi\) and \(h\) are functions of
the exit radius. By (i), \(-q>0\) on the passage, so \(\mathcal J_a\) increases;
since \(s\mathcal J_a(s\rho_{\rm g})=\mathcal J_+(\rho_{\rm g})\), we get
\(\e^{-\mathcal J_a(\varrho_e)}\le r^{c_{\rm g}/s}\). Since
\(\mathcal E\mathcal J_a=1/(\varrho_e^+K_a)\), \(\varrho_e^+\ge s\), and the Euler
derivatives of \(K_a\) are uniformly bounded, each fixed radial derivative
of \(\mathfrak b_a\) costs at most another factor \(s^{-1}\):
\(|\mathcal E^j\mathfrak b_a|\le C_js^{-j}\mathfrak b_a\), as in
\eqref{eq:direct-annulus-input-jets}. In the fixed radius box
\(\mathfrak b_a\le Cs^{-3}r^{c_{\rm g}/s-2}\); since \(r\le1/2\), this is
\(O(s^N)\) for every fixed \(N\), and in particular \(\mathfrak b_a\le1\)
for small \(s\). The chain-rule step before
\eqref{eq:direct-annulus-prefix-weight}, applied to
\(\xi=(\xi\Upsilon_a)^{-1}(\mathfrak b_a)\) with the inverse of (ii),
bounds \(|\mathcal E^j\xi|+|\mathcal E^j\log\Upsilon_a(\xi)|\) by
\(C_js^{-j}\mathfrak b_a\), which gives
\eqref{eq:positive-adaptive-input-weight}.

Logarithmic differentiation of \eqref{eq:direct-no-root-adaptive-input},
with \(\log\xi=\mathcal A_i/s+\log(h/Q)\),
\(\mathcal E\log(\varrho_e^2\varrho_e^+)=2+v_e\) and
\(\varrho_e^+(2+v_e)=2s+3\varrho_e\), gives
\[
 \mathfrak g
 =K_a^{-1}-(2s+3\varrho_e)
       -\varrho_e^+\mathcal E\log\Psi_a
       +\varrho_e^+\mathcal E\log\Upsilon_a(\xi),
\]
with \(K_a\) and \(\Psi_a(\cdot,\delta_{\rm an})\) evaluated at the exit.
By (i), the uniform Euler derivatives of \(\Psi_a\), and
\eqref{eq:positive-adaptive-input-weight}, this gives
\eqref{eq:direct-no-root-adaptive-source}.

At the cone landing, \eqref{eq:direct-adaptive-landing} gives
\(\mathcal E\log t=2+O(\epsilon_e)\) and
\(\mathcal E^2\log t=O(\epsilon_e)\). By (iii), the moving-cone
part of \cref{lem:direct-weighted-tail} applies at the landing, as before
\eqref{eq:direct-adaptive-source-slope}. With \(q_2=-s\), the endpoint
energy has derivative \((z_*/t)^{2s}a(t)\), and
\eqref{eq:direct-moving-tail} gives
\[
 |t\partial_t\log a(t)|
   \le C(\sqrt t+q_1/\sqrt t+|q_0|/t)
   \le C(\varrho_e+su_e+su_e^2).
\]
Therefore \(M_a=-\mathcal EI_+\) satisfies
\[
 M_a=t(z_*/t)^{2s}a(t)\,\mathcal E\log t,\qquad
 \mathcal E\log M_a=2-4s+
              O(\epsilon_e+\varrho_e+su_e+su_e^2).
\]
Since \(m_+=hI_{+,h}=\varrho_e^+M_a/\mathfrak g\),
\(\mathcal E\log\varrho_e^+=v_e\) and \(\partial_T=\mathcal E/\mathcal ET\),
\eqref{eq:direct-no-root-adaptive-source} gives the identity in
\eqref{eq:direct-no-root-adaptive-rate}. The factor \(K_a\) and its
derivatives have the common bounds of (i), while \(\Upsilon_a\) contributes
only terms \(O(s^N)\); this gives the stated dependence of the constants.
At the exits \(u_e\le r/\rho_{\rm g}\). Hence, for small \(m_{\rm land}\)
and radius box, after reducing \(r\) and then \(s\), the bracket in
\eqref{eq:direct-no-root-adaptive-rate} exceeds one and \(\mathfrak g<2\).
\end{proof}
